\documentclass[11pt,a4paper]{article}

\usepackage[T1]{fontenc}
\usepackage[utf8]{inputenc}
\usepackage{amsmath,amssymb,amsthm}
\usepackage{mathtools}
\usepackage{hyperref}
\usepackage{doi}
\usepackage{geometry}
\usepackage{booktabs}
\usepackage{array}
\usepackage{enumitem}

\geometry{margin=2.5cm}

\hypersetup{
  colorlinks=true,
  linkcolor=black,
  citecolor=black,
  urlcolor=black
}

\theoremstyle{definition}
\newtheorem{remark}{Remark}

\newcommand{\SU}{\mathrm{SU}}
\newcommand{\Heis}{\mathrm{Heis}}
\newcommand{\su}{\mathfrak{su}}
\newcommand{\UU}{\mathrm{U}}
\newcommand{\Mp}{\mathrm{Mp}}
\newcommand{\Spin}{\mathrm{Spin}}
\newcommand{\SL}{\mathrm{SL}}
\newcommand{\SPI}{S_\Pi}
\newcommand{\Vr}{V_\rho}
\newcommand{\Spi}{S_\Pi}
\newcommand{\Sbundle}{\mathcal{S}_\Pi}
\newcommand{\Epi}{E_\Pi}
\newcommand{\Dpi}{D_{\Pi,g,A}}
\newcommand{\Cgen}{C^3_{\mathrm{gen}}}
\newcommand{\adC}{\mathrm{ad}_\mathbb{C}}

\title{The Fermionic Matter Sub-Programme\\[0.4em]
\large Presentation Note 6}

\author{J\'{e}r\^{o}me Beau\\[0.3em]
\small Independent Researcher\\
\small \texttt{jerome.beau@cosmochrony.org}}

\date{2026}

\begin{document}

  \maketitle

  \begin{abstract}
    The fermionic matter sub-programme of the Cosmochrony corpus addresses a single
    central question: where do fermions, chirality, hypercharge, and three generations come
    from?
    The answer is that they are not postulated as external fields: they arise as the spinorial
    face of the admissible Weil module $\Vr$, once its metaplectic Lie-algebraic structure is
    complexified by the Lorentzian metric established in the geometric branch.
    The structural chain is:
    \[
      \Pi \Rightarrow F_n \simeq \Vr \Rightarrow \Mp(2,\mathbb{R})
      \Rightarrow \mathrm{mp}(2,\mathbb{R})_\mathbb{C} \simeq \mathfrak{sl}_2(\mathbb{C})
      \rightsquigarrow \Spin(3,1) \simeq \SL(2,\mathbb{C})
      \Rightarrow \Sbundle.
    \]
    Three modular structural results are established in Q14: (A)~the admissible spinor bundle
    $\Sbundle$ reproduces the $\SU(2)_L \times \UU(1)_Y$ electroweak bundle structure from
    tensor functors of $\Sbundle$; (B)~the projected Dirac operator $\Dpi$ contains a
    canonical projective endomorphism $\Epi$ that enforces the $V-A$ chiral structure
    (established via the spinorial BI lift theorem of Q14) and whose $\gamma_5$-weighted $a_4$ coefficient imposes
    anomaly-cancellation constraints making hypercharge a spectral datum; (C)~the saturation
    invariant $\sigma_c(n_3) = 3$ has a spinorial multiplicity reading yielding a gauge-singlet
    three-generation factor $\Cgen$.
    The colour-coupled quark sector is unconditional at the pointwise level
    ($[\mathrm{H\text{-}color}]_{\mathrm{pointwise}}$ established in O31).
    The structural form of $\Epi$ is now fixed as a Schur complement of the eliminated spinorial block;
    the open deliverables are its explicit Lorentzian block, the generation-split amplitude (fixed at A4
    by matching a discrete chiral-frontier normalisation with the Lorentzian Born--Infeld curvature), and
    the Yukawa sector.
  \end{abstract}

  \medskip
  \noindent\textbf{Keywords.}
  fermionic matter; admissible spinor bundle; metaplectic lift; chirality; V-A structure;
  hypercharge rigidity; anomaly cancellation; three generations; projected Dirac operator;
  Weil representation; non-injective projection; presentation note.

  \newpage
  \tableofcontents

  \newpage

  \section{Motivation and Central Question}
    \label{sec:motivation}

    The fermionic matter sub-programme answers the following question.

    \medskip
    \noindent\textit{Central question.}
    The bosonic spectral stratification — gravity at $a_2$, Yang--Mills at $a_4$ --- derives
    the bosonic sector of the Standard Model from the same admissible spectral functional.
    Fermions, chirality, hypercharge assignments, and the three-generation structure are not
    covered by the bosonic synthesis.
    \textit{Do these fermionic structures also arise as forced consequences of the admissible
    Weil fibre, or must they be introduced as additional postulates?}

    \medskip
    The answer established in Q14~\cite{Beau2026q14} is: they are forced.
    Fermions arise as the spinorial face of the Weil module $\Vr$ after Lorentzian
    complexification of its metaplectic lift.
    Chirality and the $V-A$ structure arise from the projective endomorphism $\Epi$ of the
    projected Dirac operator, via the spinorial lift of BI parity established in Q14.
    Hypercharge is not a free parameter: it is constrained by spectral anomaly cancellation.
    Three generations arise from the spinorial multiplicity reading of the same
    rank-three invariant $\sigma_c(n_3) = 3$ that yields the three spatial directions in the
    geometric branch.

    \medskip
    This sub-programme is the first paper of the Q-series that depends simultaneously on all
    four layers of the programme: the admissible Weil fibre structure (Notes~1 and~5), the
    closed Lorentzian geometry (Note~2), the gauge bundle (Note~3), and the gauge--gravity
    synthesis (Q13~\cite{Beau2026q13}, to be summarised in Presentation Note~8).

    \begin{remark}[Scope: identification only, not mass spectrum]
This sub-programme closes the \emph{identification} of the fermionic sector:
spinor bundle, chirality, hypercharge quantum numbers, and generation factor.
It does not derive the mass spectrum, the Yukawa couplings, or the CKM/PMNS mixing
matrices.
Those require the explicit construction of $\Epi$ as an operator on the spectral space,
and the integration of the projected Yukawa Lagrangian $\mathcal{L}_Y$, both of which
are listed as open deliverables.
    \end{remark}

    \begin{remark}[Vocabulary note]
Throughout this note, \emph{admissibility constraint} replaces the earlier
``relaxation'' vocabulary.
The substrate $\chi$ is static; fermions are not excitations of $\chi$ but the spinorial
interpretation of the admissible fibre after geometric complexification.
    \end{remark}

  \section{Position in the Cosmochrony Programme}
    \label{sec:position}

    The fermionic matter sub-programme sits at the apex of Branch~III.
    It depends on all five preceding sub-programmes and provides the fermionic matter
    content consumed by any downstream phenomenological application.

    Inputs from upstream:
    \begin{itemize}
      \item $F_n \simeq \Vr \simeq L^2(\mathbb{Z}/q\mathbb{Z})$, canonical symplectic structure of
      $\Vr$, and the parity involution $c \leftrightarrow q-c$ (Notes~1 and~5).
      \item Effective Lorentzian metric $g^{\mu\nu} = 2\eta^{\mu\nu}$ and effective base manifold
      $M_{\mathrm{eff}} = \mathbb{R}_\tau \times \Heis_3(\mathbb{R})$ (Note~2, Q5b--Q11).
      \item Admissible principal bundle $P_{G_\Pi}(M_{\mathrm{eff}}, G_\Pi)$ and gauge connection $A$
      (Note~3, Q6a).
      \item $\SU(3) \times \SU(2) \times \UU(1)$ gauge group (Note~3).
      \item Spectral entropy functional $\SPI[g,A]$ and the $a_4$ variation (Note~4~\cite{Beau2026sgnote}/Q12,
      Q13~\cite{Beau2026q13}).
      \item $H_{\mathrm{eff}} = \mathrm{Sym}^2(\Vr)$, three spatial directions from $\operatorname{Im}\mathbb{H}$,
      $\sigma_c(n_3) = 3$ (Note~2, Q7~\cite{Beau2026q7}).
    \end{itemize}

  \section{Logical Chain of the Sub-Programme}
    \label{sec:chain}

    The sub-programme is organised around the following structural chain.

    \begin{equation}
      \label{eq:chain}
      \Pi \Rightarrow F_n \simeq \Vr
      \;\Rightarrow\;
      \Mp(2,\mathbb{R})
      \;\Rightarrow\;
      \mathrm{mp}(2,\mathbb{R})_\mathbb{C} \simeq \mathfrak{sl}_2(\mathbb{C})
      \;\rightsquigarrow\;
      \Spin(3,1) \simeq \SL(2,\mathbb{C})
      \;\Rightarrow\;
      \Sbundle
    \end{equation}
    \[
      \;\Rightarrow\;
      \underbrace{\mathrm{Sym}^2(\Sbundle) \Rightarrow \adC(P_\Pi)}_{\SU(2)_L}
      \;,\quad
      \underbrace{\wedge^2(\Sbundle) = L_Y}_{\UU(1)_Y}
      \;,\quad
      \underbrace{\Epi \text{ left-admissible}}_{V-A}
      \;,\quad
      \underbrace{\sigma_c(n_3)=3 \to \Cgen}_{\text{3 generations}}.
    \]

    Three distinct mechanisms contribute.

    \begin{enumerate}
      \item \textit{Lorentzian complexification forces the spin group.}
      The Weil module $\Vr$ carries a canonical symplectic structure; its automorphisms lift to
      $\Mp(2,\mathbb{R})$.
      The effective metric $g^{\mu\nu} = 2\eta^{\mu\nu}$ (Lorentzian) forces the Lie algebra
      complexification $\mathrm{mp}(2,\mathbb{R})_\mathbb{C} \simeq \mathfrak{sl}_2(\mathbb{C})$,
      whose Lorentzian integration is $\Spin(3,1) \simeq \SL(2,\mathbb{C})$.
      The spinor bundle $\Sbundle$ is the associated fundamental bundle.

      \item \textit{Tensor functors of $\Sbundle$ reproduce the electroweak bundle.}
      $\mathrm{Sym}^2(\mathbb{C}^2) \simeq \mathfrak{sl}_2(\mathbb{C})$: the symmetric sector gives
      the complex adjoint, from which the compact real form selects the $\SU(2)_L$ gauge sector.
      $\wedge^2(\mathbb{C}^2) \simeq \det(\mathbb{C}^2)$: the determinant line $L_Y$ supports the
      $\UU(1)_Y$ factor.
      No independent electroweak matter bundle is introduced.

      \item \textit{The projective endomorphism $\Epi$ enforces chirality and hypercharge.}
      The projected Dirac operator satisfies the Lichnerowicz-type identity
      \[
        D^2_{\Pi,g,A}
        = -(\nabla^{S,A})^2 + \frac{R}{4} + \frac{1}{2}\gamma^\mu\gamma^\nu F^a_{\mu\nu} T^a_R
        + \Epi.
      \]
      By the spinorial lift of BI parity, established in Q14~\cite{Beau2026q14}, $\Epi$ is left-admissible
      ($P_R \Epi P_R = 0$), giving the $V-A$ structure.
      $\UU(1)_Y$ invariance of $\SPI[g,A,D_{\Pi,g,A}]$ imposes the anomaly-cancellation
      trace constraint $\mathrm{Tr}_{\Sbundle}(\gamma_5 Y A_\Pi(x)) = 0$, fixing hypercharge
      as a spectral coherence datum.
    \end{enumerate}

  \section{Constituent Papers}
    \label{sec:papers}

    The sub-programme has a single primary paper.

    \subsection{Q14: Fermionic matter and chirality from projective Dirac admissibility}
      \label{sec:q14}

      \textbf{Q14}~\cite{Beau2026q14} is the sole constituent paper.
      It is organised around three independent but mutually compatible structural results.

      \textit{Theorem~A: Spinorial electroweak bundle} (structural; unconditional).
      The admissible Weil module $\Vr$ induces, via the metaplectic lift and the Lorentzian
      complexification forced by $g^{\mu\nu} = 2\eta^{\mu\nu}$, an admissible spinor bundle
      $\Sbundle$.
      Its tensor sectors satisfy:
      \[
        \mathrm{Sym}^2(\Sbundle) \simeq \adC(P_\Pi), \qquad \wedge^2(\Sbundle) = L_Y.
      \]
      After selection of the compact real form, these define the $\SU(2)_L \times \UU(1)_Y$
      electroweak bundle data without any additional input.

      \textit{Theorem~B: Chiral selection and hypercharge rigidity}
      (structural; unconditional).
      The projected Dirac operator $\Dpi$ has the Lichnerowicz-type square above.
      By the spinorial lift of BI parity, established in Q14 (the antilinear $J_\Pi = HK$ maps
      right-chiral modes to non-admissible directions), $\Epi$ is left-admissible: $P_R \Epi P_R = 0$.
      This is the $V-A$ structure.
      $\UU(1)_Y$ invariance of the projected spectral functional then imposes, at the $a_4$ level:
      \[
        \mathrm{Tr}_{\Sbundle}(\gamma_5 Y A_\Pi(x)) = 0 \quad \forall x \in M_{\mathrm{eff}},
      \]
      where $A_\Pi$ is the $\gamma_5$-weighted $a_4$ Seeley--DeWitt density of $D^2_{\Pi,g,A}$.
      This constraint is not postulated: it is the chiral consistency condition of the same $a_4$
      level that yields Yang--Mills dynamics in the bosonic variation.
      The hypercharge weights $Y_R$ are thereby not free parameters but spectral coherence data.

      \textit{Theorem~C: Generation multiplicity} (structural; unconditional).
      The saturation invariant $\sigma_c(n_3) = 3$ has two functorially distinct readings of the
      same rank-three admissible invariant.
      In the geometric branch (Note~2), it yields three effective spatial directions.
      In the spinorial branch, the spinorial multiplicity functor maps it to a gauge-singlet
      generation space:
      \[
        \Cgen \subset \ker(\mathrm{ad}_{\SU(2)} \oplus Y).
      \]
      The identity $\mathbb{R}^3_{\mathrm{space}} \neq C^3_{\mathrm{gen}}$ is established
      explicitly (Q14 §5): the spatial and generation spaces are distinct functorial images of the
      same admissible invariant.
      Three generations are forced as a gauge-singlet factor.
      A qualitative mechanism for their splitting is identified in Q14 \S6 (summarised below); the
      explicit $\Epi$ and the amplitude of the splitting remain open.

      \textit{Dynamic generation lifting} (localisation and A4 matching).
      The internal structure of $\Cgen$ is probed by the spectrum of $\Epi^2$ (Q14 \S6).
      A static $J_\Pi$-real, weight-preserving restriction of $\Epi^2$ has a degenerate outer pair:
      the two light generations cannot be split statically (static obstruction).
      The admissible lifts removing this degeneracy form a two-dimensional $J_\Pi$-odd sector, the
      real-split direction being the weight generator $J_3$ and the mixing direction a doublet
      rotation.
      The generation split is not a finite-front observable in the symmetric sense: the raw finite front
      remains $J_\Pi$-equivariant and gives $u_{\mathrm{fin}}=0$~\cite{Beau2026prs}.
      After Lorentzian complexification, however, the chiral projector makes it meaningful to read a
      $J_\Pi$-odd frontier cocycle, and the residual automorphism is the central Weil lift
      $R_b=W(-I)=\mathrm{FT}^2$, which commutes with $\gamma_5$: it is blind to chirality and does not
      obstruct the $V-A$ selection~\cite{Beau2026q11of}.
      The remaining amplitude is fixed at A4 by matching the discrete chiral-frontier normalisation
      $\mathcal N_A$~\cite{Beau2026aar} with the Lorentzian Born--Infeld curvature
      $\mu_\chi^2=\partial_s^2\mathcal B(0)$~\cite{Beau2026prs}, neither ingredient reducing to the other.
      Generation mixing further requires the complex metaplectic phase.

      \textit{Colour sector} (unconditional at the pointwise level).
      The colour-coupled quark sector $\Sbundle \otimes V_{\mathrm{color}}$ is obtained by
      tensoring the admissible spinor bundle with the colour module $V_{\mathrm{color}}$.
      The status of this sector inherits the status of $[\mathrm{H\text{-}color}]$ from
      Notes~1 and~3: established unconditionally on the colour-adapted graph, and on the standard
      graph at the rank, averaged, and effective-exponent levels, with
      $[\mathrm{H\text{-}color}]_{\mathrm{pointwise}}$ proved analytically (O31, exact equality of
      BFS capacity profiles at finite~$q$).
      The only residual component is the extension beyond the standard Cayley graph to general
      generating sets, which is non-blocking for the quark bundle.

      \textit{Status of the spinorial BI lift.}
      The spinorial lift of BI parity to the spinor bundle $\Sbundle$ is established as a theorem
      in Q14: the BI parity involution $c \leftrightarrow q-c$, proved at the scalar Weil level in
      O18 and BornInfeld, lifts to the antilinear endomorphism $J_\Pi = HK$ satisfying
      $J_\Pi^2 = -1$ and $J_\Pi \gamma_5 = -\gamma_5 J_\Pi$, dressing the parity by the weight
      grading $H = 2J_3$.
      Chirality and $V-A$ in Q14 are therefore unconditional, their only residual dependence being
      the axiomatic chain A1--A4~\cite{Beau2026m} and the upstream BI parity datum (O30, U1).

      \begin{table}[ht]
        \centering
        \caption{Summary of Q14 results by theorem.
        Status: P = proved; S = structural.}
        \label{tab:results}
        \renewcommand{\arraystretch}{1.2}
        \begin{tabular}{@{}lllll@{}}
          \toprule
          Result & Central output & Status \\
          \midrule
          Theorem~A & Spinor bundle $\Sbundle$; EW bundle from tensor functors & S (uncond.) \\
          Theorem~B ($V-A$) & $P_R \Epi P_R = 0$; $V-A$ structure & S (uncond.) \\
          Theorem~B (hypercharge) & $Y_R$ from anomaly-cancellation trace; not free & S (uncond.) \\
          Theorem~C & $\Cgen$; three generations as gauge-singlet factor & S (uncond.) \\
          Colour sector & $\Sbundle \otimes V_{\mathrm{color}}$; quark bundle & S (uncond., pointwise) \\
          \bottomrule
        \end{tabular}
      \end{table}

    \subsection{Angular Amplitude Reduction: the $J_3$ split as oriented symplectic
    area~\cite{Beau2026aar}}
      This note fixes the numerical observable for the absolute normalisation of the angular generation
      split before any Weil--BFS computation.
      A Baker--Campbell--Hausdorff reduction identifies the per-step $J_\Pi$-odd $J_3$ component of the
      metaplectic cascade with the oriented symplectic area swept per step --- the central,
      Carnot-degree-2 increment of $\Heis_3$ --- and not with a free shear parameter.
      The quantity to measure is the accumulated oriented $J_3$-odd area along the cascade, normalised by
      the projected capacity.
      A direct evaluation on symmetric BFS shells returns $\Theta_{\mathrm{raw}}=0$ identically: the
      symmetric capacity data fix the radial sector but cannot select an oriented angular branch, so the
      orientation prescription is an irreducible input prior to the normalisation $\mathcal N_A$.

    \subsection{Q11 Oriented Frontier: recursive diagnostic and resolution of the operator
    obstruction~\cite{Beau2026q11of}}
      The symmetric-shell observable vanishes ($\Theta_{\mathrm{raw}}=0$) because the involution
      $g\mapsto g^{-1}$ pairs opposite central phases; this note defines the honest replacement, a frontier
      transfer observable on the directed edges $g\to gs$ whose central increment is the Heisenberg cocycle
      $a(g)\,s_b$, the discrete counterpart of $\tfrac12[X,Y]$ and hence of the dynamical $J_3$.
      Its main result resolves the operator obstruction: the residual reflection
      $R_b=\varphi\circ J_\Pi^{-1}=W(-I)=\mathrm{FT}^2$ is the central Weil lift, and $[\gamma_5,R_b]=0$, so
      $R_b$ is blind to chirality and the $V-A$ selector removes the $\varphi$-paired obstruction without a
      spacetime parity.
      What remains is the amplitude normalisation $\mathcal N_A$, not the operator question.

    \subsection{Projective Residue as a Schur Complement: localisation and finite/Lorentzian
    separation~\cite{Beau2026prs}}
      This note fixes the structural status of the split $u$ before the explicit $\Epi$.
      The projected Dirac square has the Schur--Feshbach form $\Epi=-M^{\dagger}M$ with
      $M=(1-P)\,D\,\Pi_S^{*}$, exhibiting $\Epi$ as the Schur complement of the eliminated spinorial block;
      $u$ is localised in the $J_\Pi$-odd part of $\Epi^2|_{\Cgen}$ as the antiunitary equivariance defect
      $\Delta_\chi(P)=\pi_{LL}-\tau\,\overline{\pi_{RR}}\,\tau^{-1}$.
      The finite locking is $J_\Pi$-equivariant, so $u_{\mathrm{fin}}=0$ (finite/Lorentzian separation): the
      split cannot originate in the finite cascade and is a Lorentzian datum, its magnitude reduced to the
      Born--Infeld saturation curvature $\mu_\chi^2=\partial_s^2\mathcal B(0)$.
      The Schur transversality on which the A4 matching was conditional is reduced, in the same note, to the
      non-vanishing of the projected commutator $[\partial_s P|_0,J_\Pi]\big|_{\mathrm{span}(e_+,e_-)}$, with an
      exhaustive transverse/null dichotomy including the zero-mode subcase.
      The determinantal Born--Infeld margin $\mathcal B(s)$, the antisymmetric structure of the chiral modulus, and the
      reduction of $\mu_\chi^2$ to the Lorentzian genus of the chiral polarisation are developed in the companion
      note~\cite{Beau2026bim}.

  \section{Inputs from Upstream Results}
    \label{sec:inputs}

    \begin{enumerate}
      \item \textbf{$F_n \simeq \Vr$ and symplectic structure} (Notes~1, 5):
      the identification of the admissible fibre as the Weil representation is the
      starting point of the metaplectic lift.
      Without this, the spinor bundle $\Sbundle$ has no canonical definition.

      \item \textbf{Lorentzian metric $g^{\mu\nu} = 2\eta^{\mu\nu}$} (Note~2, Q5b--Q11):
      the Lorentzian signature forces the complexification
      $\mathrm{mp}(2,\mathbb{R})_\mathbb{C} \simeq \mathfrak{sl}_2(\mathbb{C})$.
      Without the Lorentzian metric, the spin group would remain $\mathrm{Sp}(2,\mathbb{R})$
      and no Dirac spinors would arise.

      \item \textbf{Gauge bundle and $\SU(2)_L \times \UU(1)_Y \times \SU(3)$} (Note~3, Q6a):
      the admissible principal bundle and gauge connection $A$ enter the spin--gauge
      connection $\nabla^{S,A}$ of the projected Dirac operator.

      \item \textbf{$a_4$ Yang--Mills variation} (Q12~\cite{Beau2026q12}, Q13~\cite{Beau2026q13}):
      the spectral anomaly cancellation of Theorem~B uses the same $a_4$-level structure
      that yields Yang--Mills dynamics in the bosonic sector.
      The fermionic $a_4$ coefficient is a structural extension of the bosonic one.
      The gauge--gravity synthesis of Q13 is the scientific dependency here;
      its presentation-level treatment will appear in Presentation Note~8.

      \item \textbf{$\sigma_c(n_3) = 3$} (Note~1, O23):
      the rank-three saturation invariant provides the spinorial multiplicity input for
      Theorem~C.

      \item \textbf{BI parity involution} (BornInfeld, O18):
      the scalar BI parity $c \leftrightarrow q-c$ is the input from which the spinorial BI lift
      is derived; its derivation at the scalar level is complete.
    \end{enumerate}

  \section{Outputs to Downstream Branches}
    \label{sec:outputs}

    \begin{table}[ht]
      \centering
      \caption{Principal outputs and downstream consumers.}
      \label{tab:outputs}
      \renewcommand{\arraystretch}{1.3}
      \begin{tabular}{@{}p{5.5cm}p{4.5cm}l@{}}
        \toprule
        Output & Consumer & Status \\
        \midrule
        Admissible spinor bundle $\Sbundle$ & Phenomenology, Q-series & S \\
        $\SU(2)_L \times \UU(1)_Y$ from tensor functors & SM phenomenology & S \\
        $V-A$ chiral structure & Weak interaction pheno & S (uncond.) \\
        Hypercharge as spectral datum & SM charge assignment & S (uncond.) \\
        Three-generation factor $\Cgen$ & Mass/mixing notes & S \\
        Quark bundle $\Sbundle \otimes V_{\mathrm{color}}$ & QCD sector & S (uncond., pointwise) \\
        $\Epi$ as projective endomorphism & Yukawa sector (open) & S (open) \\
        \bottomrule
      \end{tabular}
    \end{table}

  \section{Status of Results}
    \label{sec:status}

    \subsection{Structural results (unconditional)}
      \label{sec:structural}

      \begin{enumerate}
        \item \textit{Admissible spinor bundle} (Theorem~A):
        $\Sbundle$ is induced from $\Vr$ via the metaplectic lift and Lorentzian
        complexification.
        The construction is algebraic and functorial; it depends on the Lorentzian closure
        of Q5b--Q11 but not on the spinorial BI lift.

        \item \textit{$\SU(2)_L$ from $\mathrm{Sym}^2(\Sbundle)$} (Theorem~A):
        the symmetric tensor sector gives the complex adjoint
        $\adC(P_\Pi) \simeq \mathfrak{sl}_2(\mathbb{C})$; after compact real-form selection,
        the $\SU(2)_L$ gauge sector is identified.

        \item \textit{$\UU(1)_Y$ from $\wedge^2(\Sbundle)$} (Theorem~A):
        the determinant line $L_Y = \wedge^2(\Sbundle)$ is the geometric support of the
        hypercharge factor.

        \item \textit{Generation multiplicity} (Theorem~C):
        $\Cgen \subset \ker(\mathrm{ad}_{\SU(2)} \oplus Y)$ is a gauge-singlet three-generation
        space arising from the spinorial multiplicity reading of $\sigma_c(n_3) = 3$.
        The spatial and generation readings of the rank-three invariant are distinct
        ($\mathbb{R}^3_{\mathrm{space}} \neq C^3_{\mathrm{gen}}$; Q14 §5).
      \end{enumerate}

    \subsection{Further structural results (unconditional)}
      \label{sec:further-structural}

      \begin{enumerate}
        \item \textit{Left-admissibility of $\Epi$} (Theorem~B):
        $P_R \Epi P_R = 0$ follows from the spinorial BI lift (Q14), now established.

        \item \textit{$V-A$ chiral structure} (Theorem~B):
        the condition $P_R \Epi P_R = 0$ selects the left-chiral admissible projection
        $P_L \Sbundle$, giving the $V-A$ coupling of weak interactions.

        \item \textit{Hypercharge rigidity} (Theorem~B):
        the anomaly-cancellation trace condition
        $\mathrm{Tr}_{\Sbundle}(\gamma_5 Y A_\Pi(x)) = 0$ constrains $Y_R$ to spectral
        coherence data; hypercharge is not a free parameter of the framework.

        \item \textit{Quark bundle} (§6):
        $\Sbundle \otimes V_{\mathrm{color}}$ is structural and unconditional at the pointwise
        level: the spinorial BI lift is established (Q14) and
        $[\mathrm{H\text{-}color}]_{\mathrm{pointwise}}$ is proved analytically (O31).

        \item \textit{A4 Schur-transversality reduction}~\cite{Beau2026prs}:
        the conditional residual of the A4 matching is reduced to the non-vanishing of the projected commutator
        $[\partial_s P|_0,J_\Pi]\big|_{\mathrm{span}(e_+,e_-)}$, with an exhaustive transverse/null dichotomy
        including the zero-mode subcase; the transverse branch hands the split to the Lorentzian genus, the
        Schur-null branch is marginal.
      \end{enumerate}

    \subsection{Open problems}
      \label{sec:open-problems}

      \begin{enumerate}
        \item \textit{Explicit form of $\Epi$.}
        $\Epi$ is identified as the zero-order endomorphism of the projected Dirac square, the
        ``internal spectral residue of non-injective projection.''
        Its structural form is now fixed as the Schur complement of the eliminated spinorial block,
        $\Epi=-M^{\dagger}M$ with $M=(1-P)D\,\Pi_S^{*}$~\cite{Beau2026prs}; what is not yet constructed is
        its explicit
        Lorentzian block $1-P(s)$ along the $J_\Pi$-odd modulus.
        The inter-generation splitting is localised in the $J_\Pi$-odd part of $\Epi^2$, with amplitude
        fixed at A4 by matching a discrete chiral-frontier normalisation with the Lorentzian Born--Infeld
        curvature; without the explicit Lorentzian $\Epi$, the three-generation mass spectrum and the
        Yukawa sector cannot be derived.

        \item \textit{Normalisation of $\mathcal{L}_Y$.}
        The Yukawa Lagrangian $\mathcal{L}_Y$ is identified as arising from the trace
        $\mathrm{Tr}_{\Sbundle}[\Epi^2]$ variation.
        Its integral normalisation and the derivation of the mass matrix from the spectral data
        of $\Epi$ are open.

        \item \textit{$[\mathrm{H\text{-}color}]_{\mathrm{pointwise}}$ beyond the standard graph.}
        Inherited from Notes~1 and~3.
        $[\mathrm{H\text{-}color}]_{\mathrm{pointwise}}$ is proved on the standard Cayley graph
        (O31); its extension to general generating sets is the only residual, and is non-blocking
        for the quark bundle, which is already established at the pointwise level.
      \end{enumerate}

  \section{Open Deliverables}
    \label{sec:open}

    The derivation of the spinorial BI lift, formerly the first open deliverable, is now
    established as a theorem in Q14~\cite{Beau2026q14}: the antilinear lift $J_\Pi = HK$ dresses
    the BI parity by the weight grading, making Theorems~B and~C unconditional.
    Two open deliverables remain.

    \textbf{Open deliverable 1: Explicit $\Epi$ and the Yukawa sector.}
    The mass spectrum of charged leptons and quarks, and the CKM/PMNS mixing matrices,
    require the explicit spectral construction of $\Epi$ as an operator on the admissible
    Hilbert space.
    This requires integrating the projected Yukawa Lagrangian $\mathcal{L}_Y$ and
    extracting the mass matrix from the eigenvalues of $\Epi^2$.
    This is the primary open problem of the fermionic sector.
    The structural form of $\Epi$ is now fixed as the Schur complement of the eliminated spinorial
    block~\cite{Beau2026prs};
    the open part is its explicit Lorentzian block $1-P(s)$ along the $J_\Pi$-odd modulus.
    The symmetric finite front is $J_\Pi$-equivariant and gives $u_{\mathrm{fin}}=0$, but the discrete
    chiral-frontier normalisation $\mathcal N_A$ survives the operator obstruction, since the residual
    reflection $R_b=\mathrm{FT}^2$ is blind to chirality ($[\gamma_5,R_b]=0$)~\cite{Beau2026q11of}.
    The split amplitude is therefore fixed at A4 by matching this discrete normalisation with the
    Lorentzian Born--Infeld curvature $\mu_\chi^2=\partial_s^2\mathcal B(0)$, neither ingredient reducing
    to the other:
    \[
      |u| = \mathrm{A4Match}\!\left(\mathcal N_A,\,\mu_\chi^2\right).
    \]
    The operator side of this matching is no longer conditional: by~\cite{Beau2026prs} the Schur transversality is
    reduced to the non-vanishing of the projected commutator $[\partial_s P|_0,J_\Pi]\big|_{\mathrm{span}(e_+,e_-)}$,
    with an exhaustive transverse/null dichotomy.
    Two Lorentzian data then remain open: the genus $\mathrm{sign}(\mathcal P_{\mu\nu}\mathcal P^{\mu\nu})$ in
    $g^{\mu\nu}=2\eta^{\mu\nu}$, which decides $u=0$ versus $u\neq0$ in the transverse branch~\cite{Beau2026bim}, and
    the Yukawa sector downstream of $u$.
    What remains quantitative is therefore this A4 matching, the integral normalisation of
    $\mathcal{L}_Y$, and generation mixing via the complex metaplectic phase.

    \textbf{Open deliverable 2: Full matter content in Lorentzian signature.}
    The complete fermionic content
    $S^\mathrm{matter}_\Pi = (\Sbundle \oplus (\Sbundle \otimes V_{\mathrm{color}})) \otimes \Cgen$
    in full Lorentzian signature, with explicit gauge couplings and three-generation Yukawa
    structure, is the downstream target.
    It requires the resolution of Deliverable~1 above, the colour sector being already
    established at the pointwise level (O31).

    \newpage
  \section{Conclusion}
    \label{sec:conclusion}

    The fermionic matter sub-programme establishes that fermions, chirality, hypercharge,
    and three generations are not external additions to the Cosmochrony framework.
    They are the spinorial face of the admissible Weil module $\Vr$, revealed through the
    Lorentzian complexification of its metaplectic structure.

    The electroweak bundle $\SU(2)_L \times \UU(1)_Y$ arises from the tensor functors of
    the admissible spinor bundle $\Sbundle$ (Theorem~A, unconditional).
    The $V-A$ chiral structure and hypercharge rigidity arise from the projective endomorphism
    $\Epi$ (Theorem~B, unconditional).
    Three generations arise from the spinorial multiplicity reading of $\sigma_c(n_3) = 3$
    (Theorem~C, unconditional).
    The colour-coupled quark sector follows by tensoring with $V_{\mathrm{color}}$, unconditional
    at the pointwise level (O31).

    The inter-generation splitting is localised: statically obstructed, living in a two-dimensional
    $J_\Pi$-odd lift sector, and vanishing on the symmetric finite front ($u_{\mathrm{fin}}=0$,
    $J_\Pi$-equivariant).
    Its operator obstruction is resolved, the residual reflection $R_b=\mathrm{FT}^2$ being blind to
    chirality ($[\gamma_5,R_b]=0$).
    The remaining open components are the explicit Lorentzian $\Epi$, the split amplitude fixed at A4 by
    matching the discrete chiral-frontier normalisation with the Lorentzian Born--Infeld curvature
    $\mu_\chi^2$, and the derivation of the Yukawa sector and mass spectrum.

    \bibliographystyle{plain}
    \bibliography{cosmochrony-bibliography}

\end{document}
